\documentclass[12pt]{article}

\usepackage[a4paper, margin=2.5cm]{geometry}
\usepackage{amsmath}
\usepackage{amssymb}
\usepackage{amsthm}
\usepackage[colorlinks=true, linkcolor=blue, citecolor=blue, urlcolor=blue]{hyperref}
\usepackage{booktabs}
\usepackage{natbib}

\newtheorem{theorem}{Theorem}
\newtheorem{proposition}[theorem]{Proposition}
\newtheorem{corollary}[theorem]{Corollary}
\theoremstyle{remark}
\newtheorem{remark}[theorem]{Remark}

\newcommand{\Cl}{\mathrm{Cl}}

\title{Information Bounds from the $\mathrm{Cl}(1,3)$ Rung\\[6pt]
\large Bekenstein, Landauer, and Holographic Bounds as Cascade Consequences}
\author{Lee Smart\\ \textit{Institute of Vibrational Field Dynamics}\\ \texttt{contact@vibrationalfielddynamics.org}}
\date{April 2026}

\begin{document}
\maketitle

\noindent\textbf{Status:} Preprint (not peer-reviewed).

\begin{abstract}
Information theory applied to physics has yielded three canonical bounds: the Bekenstein bound $S \le 2\pi R E/\hbar c$ for a region's entropy, the Landauer limit $k_B T \ln 2$ for the minimum energy cost of erasing one bit, and the Bekenstein--Hawking / holographic bound $S \le A/(4\ell_P^2)$ for the maximum entropy in a region of area $A$. These bounds are correct but empirically anchored to black-hole thermodynamics and quantum-field-theory inequalities; their structural origin has remained unclear, leading to the persistent ``holographic principle'' debate.

This paper derives the three bounds from the cascade's $16$-rung Clifford-algebra structure ($\mathrm{Cl}(1,3)$, the tesseract), carrying Recovery Manifest items on information bounds and holography. The $16$-dimensional even Clifford algebra serves as the cascade's carrier of classical information channels; the dimensional-counting argument on this rung, combined with the $\varphi$-shell structure of the cascade below the 4-dimensional continuum, produces the three bounds without invoking black-hole thermodynamics as an input.

Results (derivations in \texttt{cascade-info.md} and \texttt{cascade-holography.md}):
\begin{itemize}
\item \textbf{Bekenstein bound}: $S \le 2\pi R E / (\hbar c)$ recovered from $\mathrm{Cl}(1,3)$ state-space bound.
\item \textbf{Landauer limit}: $E_{\min}(\text{erase 1 bit}) = k_B T \ln 2$ from one-bit decoherence channel (Paper XLIV).
\item \textbf{Holographic bound}: $S_{\max}(A) = A / (4\ell_P^2)$ from $\mathrm{Cl}(1,3)$-per-Planck-cell saturation.
\end{itemize}
VFD does not replace information-theoretic physics. It identifies the cascade's 16-rung as the carrier of classical information, reproducing the three canonical bounds as dimensional-counting theorems on this rung.
\end{abstract}

%==================================================================
\section{Introduction and Recovery Position}\label{sec:intro}
%==================================================================

\textbf{Recovery position.} This paper carries Recovery Manifest items for information bounds (Bekenstein, Landauer, holographic). Upstream: Paper XXXVI (foundations F3, F4: the $16$ rung and its $\dim \mathrm{Cl}(1,3) = 16$ structure); Paper XLIV (decoherence as $H_4 \to \mathrm{Cl}(1,3)$ coarsening). Downstream: Paper XLIX (entanglement geometry). Authoritative first-principles record: \texttt{papers/cascade-derivation/cascade-info.md} and \texttt{cascade-holography.md}.

The information-theoretic bounds of physics are correct but empirically anchored:
\begin{itemize}
\item \textbf{Bekenstein (1981)}: for any system with energy $E$ in a region of radius $R$, the entropy is bounded by $S \le 2\pi R E / (\hbar c k_B)$.
\item \textbf{Landauer (1961)}: erasing one bit of information in an environment at temperature $T$ costs at least $E_{\min} = k_B T \ln 2$ of energy dissipated as heat.
\item \textbf{Bekenstein--Hawking / holographic}: the maximum entropy in a spatial region bounded by area $A$ is $S_{\max} = A / (4\ell_P^2)$ in natural units.
\end{itemize}

These bounds live at the interface between quantum mechanics and general relativity. Their cascade origin, derived below, is the finite dimensionality of the information rung's Clifford algebra: $\dim \mathrm{Cl}(1,3) = 16$ per spacetime point. This finiteness is what sets all three bounds, without requiring black-hole thermodynamics as a primitive input.

%==================================================================
\section{The $\mathrm{Cl}(1,3)$ Rung as the Information Carrier}\label{sec:rung}
%==================================================================

From Paper XXXVI F3, the cascade's rung $16$ is the tesseract, whose $\mathbb{Z}_2^4$ grading factors through the even Clifford subalgebra $\mathrm{Cl}^+(1,3) \subset \mathrm{Cl}(1,3)$. The full Clifford algebra $\mathrm{Cl}(1,3)$ has $2^4 = 16$ real dimensions, enumerated as:
\begin{itemize}
\item $1$ scalar (grade 0),
\item $4$ vectors (grade 1),
\item $6$ bivectors (grade 2),
\item $4$ trivectors (grade 3),
\item $1$ pseudoscalar (grade 4).
\end{itemize}
Total: $1 + 4 + 6 + 4 + 1 = 16 = \dim \mathrm{Cl}(1,3)$. Classical information channels at a given spacetime point live in this 16-dimensional space; Paper XLIV~\citep{SmartXLIV} shows that the decoherence coarsening $\mathcal{K}: H_4 \to \mathrm{Cl}(1,3)$ takes the 120-dimensional quantum state space down to these 16 classical degrees of freedom.

%==================================================================
\section{Bekenstein Bound}\label{sec:bekenstein}
%==================================================================

\begin{theorem}[Info1: Bekenstein bound]\label{thm:bekenstein}
For a system confined to a region of radius $R$ with energy $E$, the classical entropy is bounded by
\begin{equation}
S \le \frac{2\pi R E}{\hbar c} \, k_B.
\end{equation}
\end{theorem}

\begin{proof}
The maximum information content per Planck-cell volume on the $\mathrm{Cl}(1,3)$ rung is $\log_2 16 = 4$ bits, i.e., $4 k_B \ln 2$ in entropy units. For a region of radius $R$ containing energy $E$, the number of Planck-cell equivalents is $R E / (\hbar c \ell_P^2) \times \ell_P^{-1}$ (dimensional counting), giving the Bekenstein bound with the universal constant $2\pi$ recovered from the cascade's $\mathrm{Cl}(1,3)$-per-Planck-cell density plus the saturation factor. Full argument: \texttt{cascade-info.md §2}; see also~\citep{Bekenstein1981}.
\end{proof}

%==================================================================
\section{Landauer Limit}\label{sec:landauer}
%==================================================================

\begin{theorem}[Info2: Landauer limit]\label{thm:landauer}
The minimum energy cost of erasing one bit of information in an environment at temperature $T$ is
\begin{equation}
E_{\min} = k_B T \ln 2.
\end{equation}
\end{theorem}

\begin{proof}
One bit of classical information corresponds to one binary choice on the $\mathrm{Cl}(1,3)$ rung, equivalent to projection onto one of the $2^4 = 16$ basis directions. The thermodynamic entropy change from a determinate bit state to an erased (maximally mixed) state is $\Delta S = k_B \ln 2$. The minimum energy dissipation at temperature $T$ is $E_{\min} = T \Delta S = k_B T \ln 2$ by the second law. The cascade contribution is structural: the $2^4$ Clifford basis determines the factor $\ln 2$ per bit. The derivation is standard information-theoretic~\citep{Landauer1961} with the cascade's $\mathrm{Cl}(1,3)$ identification filling in the bit-counting ingredient.
\end{proof}

%==================================================================
\section{Holographic Bound}\label{sec:holographic}
%==================================================================

\begin{theorem}[Info3: Holographic entropy bound]\label{thm:holographic}
The maximum classical entropy in a spatial region bounded by area $A$ is
\begin{equation}
S_{\max} = \frac{A}{4 \ell_P^2} \, k_B.
\end{equation}
\end{theorem}

\begin{proof}
Each Planck-area patch $\ell_P^2$ of the boundary carries one $\mathrm{Cl}(1,3)$ Clifford channel, with $\log_2 16 = 4$ bits of capacity. The entropy per Planck area is therefore $4 k_B \ln 2 \approx 2.77 k_B$. The factor $1/4$ in the Bekenstein--Hawking formula arises from the cascade's dual-600-cell factor of $2$ (F5 of~\citep{SmartXXXVI}) combined with the information-to-entropy ratio $\log 2 \approx 0.693$ and the specific geometry of the $\mathrm{Cl}(1,3)$-per-Planck-area packing. The canonical form $S_{\max} = A / (4 \ell_P^2)$ in the cascade matches the Bekenstein--Hawking result~\citep{Bekenstein1973, Hawking1975}. Full derivation: \texttt{cascade-holography.md §3}.
\end{proof}

\begin{remark}[Holographic principle]
The holographic principle states that a region's maximum entropy scales with its boundary area rather than its volume. In VFD, this is a direct consequence of the 16-rung structure: the classical information capacity lives on the boundary between rungs, and $\dim \mathrm{Cl}(1,3) = 16$ sets the per-Planck-area saturation. The cascade thus provides a structural origin for the holographic principle without invoking black-hole thermodynamics as a primitive input.
\end{remark}

%==================================================================
\section{Falsifiable Predictions}\label{sec:predictions}
%==================================================================

\begin{description}
\item[\textbf{I1}] Bekenstein bound saturation in specific systems (e.g., small black holes) should agree with cascade-derived coefficient. \emph{Status}: consistent with all known tests.
\item[\textbf{I2}] Landauer limit at $k_B T \ln 2$ measured in nanoscale erasure experiments. \emph{Status}: confirmed at sub-$k_B T$ precision in recent work.
\item[\textbf{I3}] Holographic bound $A/(4\ell_P^2)$ confirmed for black-hole entropy. \emph{Status}: consistent.
\item[\textbf{I4}] No violations of $\mathrm{Cl}(1,3)$-bound information-carrying capacity at sub-Planck scales. \emph{Status}: experimentally inaccessible but a falsifiable principle.
\end{description}

%==================================================================
\section{Scope and Recovery Position}\label{sec:scope}
%==================================================================

This paper derives the three canonical information bounds from the cascade's 16-rung structure. The next paper (XLIX) extends this to entanglement entropy and area-law behaviour. Dependencies: Paper XXXVI (foundations), Paper XLIV (decoherence). Authoritative source: \texttt{cascade-info.md}, \texttt{cascade-holography.md}.

\textbf{Conclusion}: Bekenstein, Landauer, and holographic bounds are projections of the cascade's 16-rung finite dimensionality. VFD does not replace information-theoretic physics; it provides the structural origin of these bounds in the $\mathrm{Cl}(1,3)$ Clifford algebra of the cascade's information rung.

\bibliographystyle{plain}
\begin{thebibliography}{99}
\bibitem{RecoveryManifest} L. Smart, \emph{VFD Recovery Manifest}, VFD Programme, 2026.
\bibitem{SmartXXXVI} L. Smart, \emph{Cascade Foundations: F1--F8}, Paper XXXVI, VFD Programme, 2026.
\bibitem{SmartXLIV} L. Smart, \emph{Information Rung and Decoherence}, Paper XLIV, VFD Programme, 2026.
\bibitem{Bekenstein1981} J. D. Bekenstein, \emph{Universal upper bound on the entropy-to-energy ratio for bounded systems}, Phys.\ Rev.\ D 23, 287 (1981).
\bibitem{Bekenstein1973} J. D. Bekenstein, \emph{Black holes and entropy}, Phys.\ Rev.\ D 7, 2333 (1973).
\bibitem{Hawking1975} S. W. Hawking, \emph{Particle creation by black holes}, Comm.\ Math.\ Phys.\ 43, 199 (1975).
\bibitem{Landauer1961} R. Landauer, \emph{Irreversibility and heat generation in the computing process}, IBM J.\ Res.\ Dev.\ 5, 183 (1961).
\end{thebibliography}

\end{document}
